% College Lecture 6: groups, rings, fields, GF(p), GF(2^n). All polynomial results verified by script.
\section{Algebraic Structures: Groups, Rings, Fields}\label{sec:algebra}

\subsection{Groups}
A \textbf{group} $G = \langle S, \bullet \rangle$: a set with one operation satisfying
\textbf{closure}, \textbf{associativity}, \textbf{identity} $e$ ($a \bullet e = e \bullet a = a$), \textbf{inverse} ($a \bullet a' = e$). Add \textbf{commutativity} $\Rightarrow$ \textbf{abelian} (commutative) group.
\begin{itemize}
  \item $\langle \mathbb{Z}_n, + \rangle$: abelian, identity 0, inverse of $a$ is $n - a$ (so subtraction works).
  \item $\langle \mathbb{Z}_n^*, \times \rangle$: abelian, identity 1, inverses via extended Euclid (so division works).
  \item \textbf{Permutation group} (all 6 permutations of 3 elements under composition): a group but \textbf{not abelian}. Closure means two successive permutations equal one permutation, so \textbf{applying two transposition ciphers in a row is no stronger than one}.
  \item From an operation table: the identity is the element whose row copies the header; the inverse of $a$ is wherever the identity appears in $a$'s row. (Slide table on $\{a,b,c,d\}$: identity $a$, inverse pairs $(a,a), (b,d), (c,c)$.)
\end{itemize}
\textbf{Finite group}, \textbf{order} $|G|$ (number of elements). \textbf{Subgroup}: a subset that is a group under the \emph{same} operation. ($\langle \mathbb{Z}_{10}, + \rangle$ is \textbf{not} a subgroup of $\langle \mathbb{Z}_{12}, + \rangle$: the operations differ, mod 10 vs mod 12.)

\textbf{Cyclic subgroup} generated by $a$: $\{a^0, a^1, a^2, \dots\}$ (powers mean repeated application of the operation). \textbf{Cyclic group}: generated by one element, called a \textbf{generator}.
\begin{center}\small
\begin{tabular}{@{}lll@{}}
\toprule
Group & Cyclic subgroups & Generators\\
\midrule
$\langle \mathbb{Z}_6, + \rangle$ & $\{0\}$, $\{0, 2, 4\}$ (from 2, 4), $\{0, 3\}$, all of $\mathbb{Z}_6$ (from 1, 5) & 1, 5\\
$\langle \mathbb{Z}_{10}^*, \times \rangle = \{1, 3, 7, 9\}$ & $\{1\}$, $\{1, 9\}$, all of $G$ (from 3, 7) & 3, 7\\
\bottomrule
\end{tabular}
\end{center}
\textbf{Lagrange}: $|H|$ divides $|G|$. \textbf{Order of an element} $=$ order of the cyclic subgroup it generates: in $\mathbb{Z}_6$, ord$(0..5) = 1, 6, 3, 2, 3, 6$; in $\mathbb{Z}_{10}^*$, ord$(1, 3, 7, 9) = 1, 4, 4, 2$.

\subsection{Rings and fields}
\begin{center}\small
\begin{tabularx}{\textwidth}{@{}l X X l@{}}
\toprule
Structure & First operation ($+$) & Second operation ($\times$) & Example\\
\midrule
Group & all 4 (5 if abelian) & -- & $\mathbb{Z}_n$, $\mathbb{Z}_n^*$\\
Ring & all 5 & closure, associativity (+ commutativity: commutative ring); distributes over $+$ & $\langle \mathbb{Z}, +, \times \rangle$\\
Field & all 5 & all 5, except 0 has no multiplicative inverse & $\mathbb{Z}_p$, GF$(2^n)$\\
\bottomrule
\end{tabularx}
\end{center}
$\langle \mathbb{Z}, +, \times\rangle$ is a commutative ring, not a field: $12 / 5 = 2.4 \notin \mathbb{Z}$. A field supports $+, -, \times, \div$ (except $\div 0$).

\subsection{Finite fields}
\begin{keybox}
\textbf{Galois}: a finite field must have $p^n$ elements ($p$ prime, $n \ge 1$): GF$(p^n)$. For $n = 1$, GF$(p) = \mathbb{Z}_p$ with $+$ and $\times$ mod $p$.
\end{keybox}
\textbf{GF(2)} $= \{0, 1\}$: addition $=$ subtraction $=$ \textbf{XOR}, multiplication $=$ \textbf{AND}. \textbf{GF(5)}: multiplicative inverses $1 \leftrightarrow 1$, $2 \leftrightarrow 3$, $4 \leftrightarrow 4$; additive inverses $(1,4), (2,3)$.

\textbf{Why GF$(2^n)$?} Computers store $n$-bit words ($0..2^n - 1$), but $2^n$ isn't prime, so $\mathbb{Z}_{2^n}$ isn't a field. Using the largest prime below $2^n$ wastes values (for $n = 8$: 251 is the largest prime, so 251--255 are unusable). GF$(2^n)$ uses \emph{all} $n$-bit words.

\subsection{GF$(2^n)$ arithmetic}
An $n$-bit word is a polynomial of degree $\le n - 1$ with coefficients in GF(2): \ttx{10011001} $= x^7 + x^4 + x^3 + 1$; \;$x^5 + x^2 + x \to$ \ttx{00100110}.
\begin{itemize}
  \item \textbf{Addition $=$ subtraction $=$ XOR} of the bit strings (no reduction needed): $(x^5 + x^2 + x) + (x^3 + x^2 + 1)$ $=$ \ttx{00100110} $\oplus$ \ttx{00001101} $=$ \ttx{00101011} $= x^5 + x^3 + x + 1$.
  \item \textbf{Multiplication}: multiply normally (coefficients mod 2), then reduce modulo an \textbf{irreducible polynomial} of degree $n$ (cannot be factored into lower-degree polynomials; plays the role of a prime). AES uses $m(x) = x^8 + x^4 + x^3 + x + 1$ (\ttx{0x11B}).
\end{itemize}
Irreducible polynomials: degree 1: $x$, $x + 1$; degree 2: $x^2 + x + 1$; degree 3: $x^3 + x + 1$, $x^3 + x^2 + 1$; degree 4: $x^4 + x + 1$, $x^4 + x^3 + 1$, $x^4 + x^3 + x^2 + x + 1$.

\begin{keybox}[title={Worked: $(x^5 + x^2 + x) \otimes (x^7 + x^4 + x^3 + x^2 + x)$ in GF$(2^8)$}]
Product: $x^{12} + x^9 + x^8 + x^7 + x^6 + \dots$; after cancelling pairs, $x^{12} + x^7 + x^2$.\\
Divide by $x^8 + x^4 + x^3 + x + 1$: quotient $x^4 + 1$, \textbf{remainder $x^5 + x^3 + x^2 + x + 1$} (the answer).\\
\textbf{Shortcut} (repeated multiply-by-$x$): $x \cdot f$ is a left shift; if a bit falls off the top (degree 8), XOR with \ttx{00011011}. With $P_2 = x^7 + x^4 + x^3 + x^2 + x$:
\begin{center}\small
\begin{tabular}{@{}lll@{}}
$x^0 P_2$ & $x^7 + x^4 + x^3 + x^2 + x$ & \\
$x^1 P_2$ & $x^5 + x^2 + x + 1$ & reduced\\
$x^2 P_2$ & $x^6 + x^3 + x^2 + x$ & \\
$x^3 P_2$ & $x^7 + x^4 + x^3 + x^2$ & \\
$x^4 P_2$ & $x^5 + x + 1$ & reduced\\
$x^5 P_2$ & $x^6 + x^2 + x$ & \\
\end{tabular}
\end{center}
$P_1 \otimes P_2 = x^5P_2 + x^2P_2 + xP_2 = (x^6 + x^2 + x) + (x^6 + x^3 + x^2 + x) + (x^5 + x^2 + x + 1) = x^5 + x^3 + x^2 + x + 1$ \checkmark.
\end{keybox}

\textbf{Inverses}: extended Euclid on polynomials (subtraction $=$ XOR), $r_1 = m(x)$, $r_2 = a(x)$, track $t$.
\begin{center}\small
\begin{tabular}{@{}llllll@{}}
\multicolumn{6}{c}{$(x^2 + 1)^{-1}$ in GF$(2^4)$ mod $x^4 + x + 1$}\\
\toprule
$q$ & $r_1$ & $r_2$ & $r$ & $t_1$, $t_2$ & $t$\\
\midrule
$x^2 + 1$ & $x^4 + x + 1$ & $x^2 + 1$ & $x$ & $0$, $1$ & $x^2 + 1$\\
$x$ & $x^2 + 1$ & $x$ & $1$ & $1$, $x^2 + 1$ & $x^3 + x + 1$\\
$x$ & $x$ & $1$ & $0$ & $x^2 + 1$, $x^3 + x + 1$ & \\
\bottomrule
\end{tabular}
\end{center}
$\Rightarrow (x^2 + 1)^{-1} = \mathbf{x^3 + x + 1}$. Likewise $(x^5)^{-1} = x^5 + x^4 + x^3 + x$ in GF$(2^8)$ (4 rows; quotients $x^3$, $x + 1$, $x$, $x^3 + x^2 + 1$).

\textbf{GF$(2^2)$} with modulus $x^2 + x + 1$ (elements 00, 01, 10, 11 $= 0, 1, x, x+1$): $x \cdot x = x + 1$, \;$x(x + 1) = 1$, \;$(x+1)^2 = x$. GF$(2^3)$ uses $x^3 + x^2 + 1$ on the slides.

\begin{examtip}
\begin{itemize}
  \item Define group/ring/field by listing the properties per operation; give $\mathbb{Z}_n$, $\mathbb{Z}$, $\mathbb{Z}_p$ as the examples.
  \item In GF$(2^n)$: add with XOR; multiply then reduce; check degree $< n$; the past papers ask for an inverse with the full extended-Euclid table (Section~\ref{sec:pastpapers}).
\end{itemize}
\end{examtip}

\begin{fixbox}
\begin{itemize}
  \item Slide: a field's second operation ``satisfies all five properties \dots\ except that the identity of the first operation has no inverse''. Read as: every element except 0 has a multiplicative inverse.
\end{itemize}
\end{fixbox}
